\documentclass[a4paper]{article}

\usepackage{graphicx}
\usepackage{amsmath,amssymb}
\usepackage{minted}
\usepackage{multirow}
\usepackage{float}
\usepackage{todonotes}
\usepackage{forest}
\usepackage{xstring}
\usepackage{xcolor}
\usepackage[most]{tcolorbox}
\usepackage{xparse}
\usepackage{natbib}

\usetikzlibrary{graphs,graphdrawing}
\usegdlibrary{layered}

% for external graphviz
\input{/home/donkarlo/Dropbox/repo/nd_language_ai_project/src/nd_language_ai/formal/computer/programming/tex/latex/code_clips/external_graphviz.tex}

% for tags
\input{/home/donkarlo/Dropbox/repo/nd_language_ai_project/src/nd_language_ai/formal/computer/programming/tex/latex/parts/control_sequence/kind/semtag/semtag.tex}

% for links
\input{/home/donkarlo/Dropbox/repo/nd_language_ai_project/src/nd_language_ai/formal/computer/programming/tex/latex/code_clips/seehere.tex}

\title{Dependency grammar}
\date{}
\author{Mohammad Rahmani}

\begin{document}
\maketitle
    Dependency grammar (DG) is a class of modern grammatical theories that are all based on the dependency relation (as opposed to the constituency relation of phrase structure) and that can be traced back primarily to the work of Lucien Tesnière. Dependency is the notion that linguistic units, e.g. words, are connected to each other by directed links. The (finite) verb is taken to be the structural center of clause structure. All other syntactic units (words) are either directly or indirectly connected to the verb in terms of the directed links, which are called dependencies \seeherefooter{https://en.wikipedia.org/wiki/Dependency_grammar}.
    \bibliographystyle{plainnat}
    \bibliography{/home/donkarlo/Dropbox/repo/nd_shared_research/bibliography/bibliography.bib}
\end{document}